\chapter{Synovial Fluid Analysis}
\section*{What Is This Test?}
Synovial fluid analysis evaluates fluid aspirated from a joint space (arthrocentesis). The analysis includes: macroscopic appearance (color, clarity, viscosity), total white blood cell (WBC) count with differential, Gram stain and culture, crystal analysis by compensated polarized light microscopy, and glucose and protein. Synovial fluid analysis is essential for the diagnosis of septic arthritis (a rheumatologic emergency requiring immediate IV antibiotics and surgical drainage/washout), crystal-induced arthritis (gout --- monosodium urate [MSU] crystals; pseudogout/calcium pyrophosphate deposition disease [CPPD] --- calcium pyrophosphate dihydrate [CPPD] crystals), and inflammatory versus non-inflammatory arthritis. Joint fluid is classified as: non-inflammatory (WBC <2,000 cells/mcL, <25\% PMNs, e.g., osteoarthritis, trauma), inflammatory (WBC 2,000--50,000, >50--75\% PMNs, e.g., rheumatoid arthritis, gout, pseudogout, reactive arthritis, psoriatic arthritis, SLE), and septic (WBC >50,000--100,000+, >90\% PMNs, Gram stain positive, culture positive, markedly low glucose [<50\% of serum]).
\section*{Alternative Names}
Joint Fluid Analysis; Arthrocentesis Fluid; Synovial Fluid Cell Count and Differential; Crystal Analysis
\section*{Principle}
Synovial fluid is collected by arthrocentesis under sterile conditions. WBC count determined by hemocytometer. Differential performed on cytocentrifuged Wright-Giemsa-stained preparation. Crystal analysis: a drop of fresh fluid is examined under compensated polarized light microscopy with a first-order red compensator. MSU crystals (gout): needle-shaped, 2--20 mcm long, negatively birefringent (yellow when parallel to the compensator, blue when perpendicular). CPPD crystals (pseudogout): rhomboid or rectangular, positively birefringent (blue when parallel, yellow when perpendicular). Basic calcium phosphate (BCP, apatite) crystals are NOT visible by polarized light microscopy and require alizarin red S staining, electron microscopy, or spectroscopy.
\section*{History}
Arthrocentesis has been performed since the late 1800s, when physicians first aspirated joint effusions to relieve symptoms and inspect the fluid. In 1961, McCarty and Hollander introduced compensated polarized light microscopy to identify monosodium urate crystals, establishing gout as a crystal-induced arthropathy. The following year, McCarty and colleagues described calcium pyrophosphate dihydrate crystals and coined the term ``pseudogout syndrome,'' now termed calcium pyrophosphate deposition disease. In the same era, investigators defined the synovial fluid white blood cell thresholds that distinguish non-inflammatory, inflammatory, and septic effusions, cementing arthrocentesis as a diagnostic cornerstone of the acute monoarthritis evaluation. Standardized guidelines for the evaluation and management of the hot swollen joint followed in the 1970s and 1980s.
\section*{Why This Test Is Ordered}
\begin{itemize}
  \item Rule out septic arthritis in any acutely swollen, painful joint: a medical emergency requiring urgent arthrocentesis, Gram stain, and culture, because infection can destroy a joint within 24--48 hours
  \item Distinguish crystal-induced arthritis: gout (monosodium urate crystals) versus pseudogout (calcium pyrophosphate dihydrate crystals)
  \item Classify an effusion as non-inflammatory, inflammatory, or septic to narrow the differential diagnosis of monoarthritis and polyarthritis
  \item Confirm the diagnosis of gout or pseudogout by demonstrating crystals, which is the only definitive method; a clinical diagnosis is presumptive
  \item Evaluate chronic monoarthritis when imaging and blood tests are nondiagnostic
  \item Confirm resolution of infection and monitor response in patients treated for septic arthritis
\end{itemize}
\section*{Is This a Primary or Secondary Test}
Synovial fluid analysis is the primary diagnostic test for septic arthritis and crystal-induced arthritis; joint aspiration and fluid examination are the definitive studies when these diagnoses are suspected. For systemic inflammatory arthritis (rheumatoid arthritis, reactive arthritis, psoriatic arthritis, SLE), it is a secondary, supportive test used to exclude infection and crystals before the diagnosis is established.
\section*{How This Test Is Performed}
Arthrocentesis is performed under strict sterile technique after skin preparation and local anesthesia. Fluid is aspirated from the largest accessible effusion, most often the knee, and immediately distributed to a plain tube for cell count, a heparinized tube for crystal analysis, and sterile culture bottles. A drop of fresh fluid is placed on a slide for Gram stain, and a separate fresh drop is examined under compensated polarized light promptly because crystals dissolve and cells degrade with storage. The specimen is assessed for volume, color, clarity, and viscosity, and a string test is performed to gauge mucin content.
\section*{What Preparation Is Needed}
No fasting or dietary preparation is required. Anticoagulant and antiplatelet medications are usually continued but should be disclosed because they increase bleeding risk; when feasible, a platelet count and INR are obtained before aspiration in patients on anticoagulation. Topical anesthetic may be applied before the procedure.
\section*{Contraindications and Precautions}
Contraindications: overlying skin infection (cellulitis) at the puncture site, bacteremia (relative), and severe uncorrected coagulopathy (relative). Cautions: the fluid must be analyzed promptly because white cell counts decline and crystals dissolve or form with storage; liquid EDTA anticoagulant can dissolve CPPD crystals, so plain or heparinized tubes are used for crystal analysis; and samples obtained before antibiotics are essential for culture yield.
\section*{Risks}
Low risk. Pain or bruising at the puncture site, bleeding into the joint (hemarthrosis), infection introduced by aspiration (rare, less than 1 in 10,000 for native joints), and cartilage damage if the needle scrapes bone. Vasovagal reactions may occur.
\section*{What Happens After the Test}
Cell count, crystal examination, and Gram stain results are typically available within hours; culture results require 48--72 hours. If septic arthritis is suspected, empirical intravenous antibiotics are started immediately and are never withheld pending culture results, and urgent surgical drainage may be required. The joint is rested and the patient observed for increasing pain or swelling that suggests a complication.
\section*{Understanding Results}
\subsection*{Non-Inflammatory (WBC <2,000/mcL)}
Osteoarthritis, traumatic effusion, hemarthrosis, avascular necrosis, and internal derangement (meniscal/ligamentous injury).
\subsection*{Inflammatory (WBC 2,000--50,000/mcL)}
Rheumatoid arthritis, gout, pseudogout, reactive arthritis, psoriatic arthritis, systemic lupus erythematosus, and ankylosing spondylitis.
\subsection*{Septic (WBC >50,000/mcL, often >100,000)}
Bacterial septic arthritis (S. aureus, Streptococcus species, N. gonorrhoeae in young sexually active adults). Gram stain positive in 50--75\%. Culture positive in 90--95\% if no prior antibiotics. Glucose markedly decreased (<50\% of serum glucose). Lactate elevated. Medical emergency: risk of rapid irreversible joint destruction within 24--48 hours.
\section*{Normal Results}
Volume: <3.5 mL (knee). Color: colorless to pale yellow. Clarity: clear. Viscosity: high (string test 3--5 cm). WBC: <200/mcL. Differential: <25\% PMNs. No crystals. Gram stain and culture: negative. Glucose: approximately equal to serum glucose. Protein: 1.0--3.0 g/dL.
\section*{Follow-up Tests}
Blood cultures (before antibiotics if septic arthritis suspected), serum uric acid, serum rheumatoid factor and anti-CCP, synovial fluid culture (aerobic, anaerobic, fungal, mycobacterial as indicated), and synovial biopsy if chronic monoarthritis with negative cultures.
\section*{References}
\begin{enumerate}
  \item Margaretten ME, Kohlwes J, Moore D, Bent S. Does this adult patient have septic arthritis? JAMA. 2007;297(13):1478--1488.
  \item Pascual E, Jovani V. Synovial fluid analysis. Best Pract Res Clin Rheumatol. 2005;19(3):371--386.
  \item Shmerling RH. Synovial fluid analysis. Rheum Dis Clin North Am. 1994;20(2):503--512.
  \item McCarty DJ, Hollander JL. Identification of urate crystals in gouty synovial fluid. Ann Intern Med. 1961;54:452--460.
  \item Swan A, Amer H, Dieppe P. The value of synovial fluid assays in the diagnosis of joint disease: a literature survey. Ann Rheum Dis. 2002;61(6):493--498.
\end{enumerate}
\clearpage
